% sec04c-ab.tex — Section 4 subsection: the (b,beta_r)/(beta_r,b) channels
% Writer: 作家_AB (Writer_AB). Body LaTeX only; macros live in main.tex.

\subsection{The channels $(b,\beta_r)$ and $(\beta_r,b)$}\label{sec:superspace:ab}

We now compute the mixed family $\nabla_-(b^{A}\beta_r^{B})$ and $\nabla_-(\beta_r^{A}b^{B})$ in full, taking $r=1$ without loss (flavor transport gives $r=2,3$). It is the only family whose output word is not read off a cut parent graph directly: its raw graph carrier contains an equation-of-motion (total-derivative) layer that must be quotiented, and a finite normal-product term that must be added, before the renormalized answer lands on the channel-table row~\eqref{eq:ss:ch-bbeta}. Graph weights use unit-normalized insertions $\nabla_+\cW_+$, $\nabla_+\Phi_r$; the letter normalizations~\eqref{eq:ss:letters} are restored at the end.

\paragraph{1. Graph inventory.} Three ordered-triangle families contribute, with routing, external slot map, and denominators
\begin{equation}\label{eq:ss:ab:routing}
r_0=\ell,\qquad r_1=\ell-p,\qquad r_2=\ell-p-q,\qquad
\partial c=\partial c\,(p),\qquad \beta_r=\beta_r\,(q),\qquad D_i=r_{i,d}^{2}
\end{equation}
(the reflected slot map $\beta_r(p),\partial c(q)$ is false): \textbf{G1}, the gauge triangle (three vector internal lines, one matter-cubic and one gauge-cubic vertex; the $\nabla_+\cW_+$ mark tags $r_2$, the $\nabla_+\Phi_r$ mark tags $r_0$), output carrier $\langle\partial c,\beta_r\rangle$ for $(b,\beta_r)$; \textbf{G2}, the mixed triangle (one vector and two matter internal lines, two matter-cubic vertices), output $\langle\beta_r,\partial c\rangle$ for $(b,\beta_r)$; \textbf{G3}, two ordered triangles $G_{3,2},G_{3,3}$ (one matter-cubic and one antichiral superpotential vertex; internal flavors $(1,2)$, $(1,3)$), outputs $\langle\gamma^2,\gamma^3\rangle$, $\langle\gamma^3,\gamma^2\rangle$. Each ordering is reflected for $(\beta_r,b)$; the line contents are displayed in Figure~\ref{fig:ss:ab:triangles}.
\begin{figure}[t]
\centering
\begin{tikzpicture}[baseline]
\node[svertex] (v0) at (90:1.5) {}; \node[svertex] (v1) at (210:1.5) {}; \node[svertex] (v2) at (330:1.5) {};
\draw[sV] (v1) -- node[midway,sinsert,label=left:{$\nabla_{+}\Phi_{r}$}] {} (v0); \draw[sV] (v0) -- (v2);
\draw[sV] (v2) -- node[midway,sinsert,label=below:{$\nabla_{+}\cW_{+}$}] {} (v1);
\draw[sleg] (v0) -- ++(90:0.8) node[above] {$\partial_{\dot\alpha}c\,(p)$};
\draw[sleg] (v2) -- ++(330:0.8) node[anchor=north west] {$\beta_{r}\,(q)$};
\node at ($(v1)!0.5!(v0)+(0.45,0)$) {$r_{0}$}; \node at ($(v0)!0.5!(v2)+(-0.45,0)$) {$r_{1}$}; \node at ($(v2)!0.5!(v1)+(0,0.4)$) {$r_{2}$};
\node at (0,-2.2) {G1: three vector internal lines};
\begin{scope}[xshift=6.5cm]
\node[svertex] (w0) at (90:1.5) {}; \node[svertex] (w1) at (210:1.5) {}; \node[svertex] (w2) at (330:1.5) {};
\draw[sPhi] (w1) -- node[midway,sinsert,label=left:{$\nabla_{+}\Phi_{r}$}] {} (w0); \draw[sV] (w0) -- (w2);
\draw[sPhit] (w2) -- node[midway,sinsert,label=below:{$\nabla_{+}\cW_{+}$}] {} (w1);
\draw[sleg] (w0) -- ++(90:0.8) node[above] {$\partial_{\dot\alpha}c\,(p)$};
\draw[sleg] (w2) -- ++(330:0.8) node[anchor=north west] {$\beta_{r}\,(q)$};
\node at ($(w1)!0.5!(w0)+(0.45,0)$) {$r_{0}$}; \node at ($(w0)!0.5!(w2)+(-0.45,0)$) {$r_{1}$}; \node at ($(w2)!0.5!(w1)+(0,0.4)$) {$r_{2}$};
\node at (0,-2.2) {G2/G3: two matter lines, one vector line};
\end{scope}
\end{tikzpicture}
\caption{The ordered triangles of the $(b,\beta_r)/(\beta_r,b)$ family in the routing~\eqref{eq:ss:ab:routing}: one $\nabla_+\cW_+$ and one $\nabla_+\Phi_r$ insertion (filled dots), residual external legs $\partial c(p)$, $\beta_r(q)$. The right panel shows the G2 line content (vector edge $r_1$, legs $\partial c,\beta_r$); the two G3 routings share the vertex structure with the vector edge on $r_0$ and the antichiral legs $\gamma^s,\gamma^t$ in place of $\partial c,\beta_r$, one matter-cubic vertex replaced by the antichiral superpotential vertex.}
\label{fig:ss:ab:triangles}
\end{figure}
Excluded from the graph sum: the zero-square current bubbles $JE,JX,PE,PX$, which cancel in exact pairs ($I_{JE}+I_{JX}=0$, $I_{PE}+I_{PX}=0$) and contain no four-dimensional inverse square, so each DRED defect is zero; and the separately drawn source-resolvent bubbles, which are presentations of the same induced derivative-of-action contact and are not independent anomaly graphs.

\paragraph{2. Wick contractions.} The Euclidean propagators in the canonically normalized frames are
\begin{equation}\label{eq:ss:ab:propagators}
\langle\cV^{A}(p,1)\cV^{B}(-p,2)\rangle_E
=-\frac{2\h g^{2}\kappa^{AB}}{p^{2}}\,\delta^{4}(\theta_{12}),
\qquad
\langle\Phi_r^{A}(p,1)\widetilde\Phi^{sB}(-p,2)\rangle_E
=\delta_r{}^{s}\,\frac{\h g^{2}\kappa^{AB}}{16\,p^{2}}\,
\bcD_1^{2}\cD_1^{2}\,\delta^{4}(\theta_{12}) .
\end{equation}
For G3 the flavor- and chirality-preserving Wick pairing is unique: the vector legs of the source and matter vertices pair along $r_0$, the matter leg of the matter vertex pairs with an antichiral leg of the superpotential vertex along $r_1$, and the source matter leg pairs with the remaining antichiral leg along $r_2$; the two uncontracted legs are the external legs. There is no second preserving pairing and no graph automorphism division, so
\begin{equation}\label{eq:ss:ab:cprop}
C_{\mathrm{prop}}=(-\h)\left(\frac{\h}{16}\right)^{2}=-\frac{\h^{3}}{256} .
\end{equation}
Every primitive weight is one:
\begin{equation}\label{eq:ss:ab:mult}
\left[\tfrac12(2)_{\text{cyc.\ source}}\right]
\left[\tfrac1{2!}(2)_{\text{vertex orders}}\right]
\left[\tfrac1{3!}(6)_{\text{gauge-vertex labels}}\right]
\left[\tfrac1{3!}(6)_{\cV\Phi\widetilde\Phi}\right]
(1)_{\mathrm{Wick}}=1\ \text{(G1)},\qquad
\left[\tfrac12(2)\right]
\left[\tfrac1{2!}(2)\right]
\left[\tfrac1{3!}(6)_{\text{superpotential labels}}\right]
(1)_{\mathrm{Wick}}=1\ \text{(G3)} .
\end{equation}

\paragraph{3. Amplitude assembly.} The mixed source Hessian, the matter-cubic vertex, and the antichiral superpotential vertex (from $W$ of Section~\ref{sec:notation}) contribute, with~\eqref{eq:ss:ab:cprop} and the unit factorial weights~\eqref{eq:ss:ab:mult},
\begin{equation}\label{eq:ss:ab:vertexconstants}
C_I=-\frac{g^{2}}{4\sqrt2},
\qquad
C_M=\frac{\sqrt2\,g}{\h},
\qquad
C_{H_-}=-\frac{\sqrt2\,g}{\h};
\qquad
C_{\mathrm{preD}}^{G2}=C_I C_M^{2}C_{\mathrm{prop}}=+\frac{\sqrt2\,\h g^{4}}{1024},
\qquad
C_{\mathrm{preD}}^{32}=C_I C_M C_{H_-}C_{\mathrm{prop}}=-\frac{\sqrt2\,\h g^{4}}{1024},
\qquad
C_{\mathrm{preD}}^{33}=+\frac{\sqrt2\,\h g^{4}}{1024},
\end{equation}
the relative sign of $G_{3,2},G_{3,3}$ being the ordered flavor/color reflection sign of the superpotential vertex. The residual external-leg operators are: for G1, the pairing $X_{\partial c\,\beta}:=\langle\partial c,\beta_r\rangle$ and the equation-of-motion composite $E:=\beta_r\,(P^{\dot\alpha}\partial_{\dot\alpha}c)$ ($P^{\dot\alpha}$ the holomorphic momentum of the $\partial c$ leg in Fourier space); for G2, the single word $\tfrac13\,\beta_r(4p-q)^{\dot\alpha}\partial_{\dot\alpha}c$; for G3, the uncontracted $\gamma^s,\gamma^t$ legs, contracted by the wedge $W_{ij}:=r_{i,+}\wedge r_{j,+}=r_{i,+\dot\alpha}r_{j,+}{}^{\dot\alpha}$, $W_{ji}=-W_{ij}$.

\paragraph{4. The $D$-algebra chain.} Discarding a longitudinal word at edge $e$ against the neighboring projector at edge $n$ transports the occurrence by superspace integration by parts; the endpoint sign, the corner identity, and the loop-saturation $16$ give the universal transport factor and its full-$d$ companion
\begin{equation}\label{eq:ss:ab:transport}
\left(-\frac12\right)(-1)(16)=8,
\qquad
8\,\bar r_n^{\,2}\,\mathcal T_{e|n}-8\,r_{n,d}^{\,2}\,\mathcal T_{e|n}
=8\,\mul^{2}\,\mathcal T_{e|n} .
\end{equation}
For G3 the raw Grassmann-measure monomial, converted to the full measure, settles the unit
\begin{equation}\label{eq:ss:ab:unit4096}
32768\left(\tfrac14\right)_{d^{4}\theta_M}
\left(\tfrac12\right)_{d^{2}\bar\theta_H}
\left(-\tfrac12\right)_{\mul^{2}\text{ extraction}}
(2)_{\text{trace/metric}}=-4096 ,
\end{equation}
the trace/metric factor $2$ being algebraic and canceled by the $-\tfrac12$ rank extraction, not a graph multiplicity. The three kinetic cuts and their parents are, with $W_0=W_{12}$, $W_1=W_{P0}$, $W_2=W_{p0}$,
\begin{align}
K_0&=+4096\,W_0,
&
K_1&=-4096\,W_1,
&
K_2&=+4096\,W_2,
\label{eq:ss:ab:cuts}\\
P_0&=-4096\,(D_0+\mul^{2})\,W_0,
&
P_1&=+4096\,(D_1+\mul^{2})\,W_1,
&
P_2&=-4096\,(D_2+\mul^{2})\,W_2 .
\label{eq:ss:ab:parents}
\end{align}
For G2 the selected and transported occurrences give
\begin{equation}\label{eq:ss:ab:g2pair}
-\frac{8\,(D_2+\mul^{2})\,\mathcal R}{D_0D_1D_2}
+\frac{8\,\mathcal R}{D_0D_1}
=-\frac{8\,\mul^{2}\,\mathcal R}{D_0D_1D_2},
\qquad
-\tfrac12\,(D_1+\mul^{2})\,W+\tfrac12\,D_1W=-\tfrac12\,\mul^{2}W ,
\end{equation}
and every G1 row obeys $F_A+C_{A,d}=-8R_A\,\mul^{2}$, $F_B+C_{B,d}=-8R_B\,\mul^{2}$, with $C_{A,d}=-L_A-\alpha_A r_{2,d}^{2}$, $\alpha_A=-8R_A$, $L_A=-\tfrac12\cD_+\bcD^{2}\cD^{2}$. The wedge algebra of the three G3 remainders closes pointwise: with $r_1=r_0-p$, $r_2=r_0-p-q$,
\begin{equation}\label{eq:ss:ab:wedge}
-W_{12}+W_{P0}-W_{p0}=-\,p_+\wedge q_+
\quad\Longrightarrow\quad
\mathcal R_{G3}^{\mathrm{full}}
=-\frac{4096\,\mul^{2}\,(p_+\wedge q_+)}{D_0D_1D_2} .
\end{equation}

\paragraph{5. Evanescent extraction and the rank-one moments.} Every remainder of step 4 is governed by the cutting-failure identity~\eqref{eq:ss:cuttingfailure} and the master integral~\eqref{eq:ss:masterintegral}. The rank-one moments are the ordered-simplex averages
\begin{equation}\label{eq:ss:ab:moments}
\lim_{\epsilon\to0}\int_{\ell}\frac{\mul^{2}\,\bar r_0^{\mu}}{D_0D_1D_2}
=\frac{2p^{\mu}+q^{\mu}}{96\pi^{2}},
\qquad
\lim_{\epsilon\to0}\int_{\ell}\frac{\mul^{2}\,\bar r_1^{\mu}}{D_0D_1D_2}
=\frac{q^{\mu}-p^{\mu}}{96\pi^{2}},
\qquad
\lim_{\epsilon\to0}\int_{\ell}\frac{\mul^{2}\,\bar r_2^{\mu}}{D_0D_1D_2}
=-\frac{p^{\mu}+2q^{\mu}}{96\pi^{2}},
\qquad
\lim_{\epsilon\to0}\int_{\ell}\frac{\mul^{2}\,W_{p0}}{D_0D_1D_2}
=\frac{W(p,q)}{96\pi^{2}},
\end{equation}
from $2\int_{\Sigma_2}x_i\,dx_1dx_2=\tfrac13$ per coordinate. Applied graph by graph: G1 integrates to the mark coefficients $(\tfrac43,\tfrac23)$, $(\tfrac23,\tfrac43)$ in the two slot orderings, i.e.\ the word $2(p+q)^{\dot\alpha}\beta_r\,\partial_{\dot\alpha}c=2\,P^{\dot\alpha}\bigl(\beta_r\,\partial_{\dot\alpha}c\bigr)$, a pure total derivative vanishing in the quotient; G2 integrates to the selected word $\tfrac23\,\beta_r(2p+q)^{\dot\alpha}\partial_{\dot\alpha}c$ plus the transported word $-\beta_r\,q^{\dot\alpha}\partial_{\dot\alpha}c$, i.e.\ $\tfrac13\,\beta_r(4p-q)^{\dot\alpha}\partial_{\dot\alpha}c$ with pair/equation-of-motion coefficients $(-\tfrac13,\tfrac43)$ on the graph-word placement (Remark~\ref{rem:ss:ab:slotmap}); G3 carries the marked simplex weights $\tfrac23,\tfrac13$ per mark and sums to the raw wedge coefficients $+2\sqrt2\,i$ ($G_{3,2}$) and $-2\sqrt2\,i$ ($G_{3,3}$) in units of $\h g^{2}/(16\pi^{2})$. Passing from the raw momentum wedge to the dotted bilinear flips the sign once, since with the locked phase $e^{ipx}$, $(ip)_+\wedge(iq)_+=-(p_+\wedge q_+)$.

\paragraph{6. Pole cancellation.} No $1/\epsilon$ pole occurs anywhere in this family: every parent-plus-contact pair vanishes identically at $\mul^{2}=0$,
\begin{align}
\frac{P_0}{D_0D_1D_2}+\frac{K_0}{D_1D_2}
&=-\frac{4096\,\mul^{2}W_0}{D_0D_1D_2},
&
\frac{P_1}{D_0D_1D_2}+\frac{K_1}{D_0D_2}
&=+\frac{4096\,\mul^{2}W_1}{D_0D_1D_2},
&
\frac{P_2}{D_0D_1D_2}+\frac{K_2}{D_0D_1}
&=-\frac{4096\,\mul^{2}W_2}{D_0D_1D_2},
\label{eq:ss:ab:polecancel}
\end{align}
and likewise for the G2 pairs~\eqref{eq:ss:ab:g2pair} and the G1 rows. Only the evanescent remainders survive. (Contrast the seed channel $(b,b)$, where literal $1/\epsilon$ poles cancel between the parent triangle and the Schwinger contact rows.)

\paragraph{7. Raw carrier, total-derivative quotient, finite shift.} Collect the integrated outputs as coefficient vectors of unit-normalized carriers in the ordered basis of candidate bilinears
\begin{equation}\label{eq:ss:ab:basis}
\bigl(X_{\partial c\,\beta},\,E,\,X_{\beta\,\partial c},\,E,\,
C_{st},\,C_{ts}\bigr),
\qquad
X_{\beta\,\partial c}:=\langle\beta_r,\partial c\rangle,
\quad
C_{st}:=\langle\gamma^{s},\gamma^{t}\rangle,
\quad
C_{ts}:=\langle\gamma^{t},\gamma^{s}\rangle,
\end{equation}
with $\varepsilon_{rst}=+1$ and $E:=\beta_r\bigl(P^{\dot\alpha}\partial_{\dot\alpha}c\bigr)$ the equation-of-motion composite of step 3, entering in both orderings. The two pairings are a single carrier: the odd--odd Koszul swap ($|\partial c|=|\beta_r|=1$) and the dotted-index flip $\psi^{\dot\alpha}\chi_{\dot\alpha}=-\psi_{\dot\alpha}\chi^{\dot\alpha}$ of the locked convention cancel,
\begin{equation}\label{eq:ss:ab:paireq}
X_{\partial c\,\beta}=(\partial_{\dot\alpha}c)(\partial^{\dot\alpha}\beta_r)
=-(\partial^{\dot\alpha}\beta_r)(\partial_{\dot\alpha}c)
=(\partial_{\dot\alpha}\beta_r)(\partial^{\dot\alpha}c)
=X_{\beta\,\partial c}=:X .
\end{equation}
There is accordingly one total-derivative identity, not two: the even derivation $P^{\dot\alpha}$ acting on the odd--odd product gives no Koszul sign, and the pair term enters with the index flip $(P^{\dot\alpha}\beta_r)\partial_{\dot\alpha}c=-\langle\beta_r,\partial c\rangle=-X$,
\begin{equation}\label{eq:ss:ab:td}
T:=P^{\dot\alpha}\bigl(\beta_r\,\partial_{\dot\alpha}c\bigr)
=\bigl(P^{\dot\alpha}\beta_r\bigr)\partial_{\dot\alpha}c
+\beta_r\bigl(P^{\dot\alpha}\partial_{\dot\alpha}c\bigr)
=-X+E ,
\end{equation}
so $aX+bE=(a+b)X+bT$ in either ordering. (Equivalently $P^{\dot\alpha}(\partial_{\dot\alpha}c\,\beta_r)=X-E=-T$: a single ray. The combination $X+E$ is not a total derivative, and $X$ itself is not one either.) The raw pair coefficients displayed below are quoted on the graph-word placement $X_{>}:=(\partial^{\dot\alpha}\beta_r)\partial_{\dot\alpha}c=-X$, where the same identity reads $T=X_{>}+E$, i.e.\ $aX_{>}+bE=(a-b)X_{>}+bT$: thus the G1 entries $(2,2)_{(X_{>},E)}=(0,2)_{(X_{>},T)}$ quotient to zero, and the G2 split and quotient are the machine-sealed replay values of Remark~\ref{rem:ss:ab:slotmap}. The coefficient flow of the family is the aligned display
\begin{align}
\text{raw graph carrier:}\quad
&\bigl(2,\;2,\;-\tfrac13,\;\tfrac43,\;-2i\sqrt2,\;+2i\sqrt2\bigr)
\quad\text{on the six-carrier basis }\eqref{eq:ss:ab:basis},
\label{eq:ss:ab:flowraw}\\[2pt]
\text{total-derivative quotient:}\quad
&\bigl(0,\;1,\;-2i\sqrt2,\;+2i\sqrt2\bigr)
\quad\text{on }
\bigl(X_{\partial c\,\beta},X_{\beta\,\partial c},C_{st},C_{ts}\bigr),
\label{eq:ss:ab:flowtd}\\[2pt]
\text{finite normal-product shift:}\quad
&+\;\bigl(1,\;0,\;+i\sqrt2,\;-i\sqrt2\bigr),
\label{eq:ss:ab:flowshift}\\[2pt]
\text{renormalized vector:}\quad
&\bigl(1,\;1,\;-i\sqrt2,\;+i\sqrt2\bigr) .
\label{eq:ss:ab:flowren}
\end{align}
The shift is fixed by residual supersymmetry: acting with the $q_r$ of~\eqref{eq:ss:qaction} (stated with the channel table in Section~\ref{sec:superspace:table}) on a general vector $x=(x_1,x_2,x_3,x_4)$ in the four-component basis, the Ward conditions of the descendant system $q_s[\,x\cdot(\text{basis})\,]=i\delta_{sr}\,\Delta(b,b)$ of~\eqref{eq:ss:triplet} --- the homogeneous $s=2,3$ mixing plus the $s=1$ matching against the anchored $(b,b)$ singlet, restricted to this flavor ansatz --- read
\begin{equation}\label{eq:ss:ab:ward}
x_1-x_2=0,
\qquad
i\sqrt2\,x_1+x_3=0,
\qquad
-i\sqrt2\,x_1+x_4=0 ,
\end{equation}
a rank-three system whose solution space is the ray $x=t\,(1,1,-i\sqrt2,+i\sqrt2)$: the renormalized vector~\eqref{eq:ss:ab:flowren} is exactly annihilated by~\eqref{eq:ss:ab:ward}, and the descendant identity~\eqref{eq:ss:triplet} anchors the ray to the $(b,b)$ scale, fixing $t=1$.

\begin{remark}[The finite shift is a scheme term, not a new anomaly graph]%
\label{rem:ss:ab:shift}
The vector~\eqref{eq:ss:ab:flowshift} is a finite composite-source normal-product (scheme) term. It is not a triangle, collapsed graph, Schwinger cut, inverse-square occurrence, or new cutting failure: all bare-graph anomaly of this family enters through the $\mul^{2}$ remainders of steps 4--6, and the shift is precisely what removes the residual factor of two between the raw G3 $\gamma\gamma$ coefficients and the Ward-compatible ray.
\end{remark}

\begin{remark}[Granularity of the machine replays]\label{rem:ss:ab:machine}
Three ingredients are sealed at machine level and displayed here at the granularity of the extracted formulas: the full 192-row G1 $D$-word table (only test-point rows admit an analytic display; the row identities $F_A+C_{A,d}=-8R_A\mul^{2}$, $F_B+C_{B,d}=-8R_B\mul^{2}$ hold for every row), the replay producing the G2 raw residual, and the color-route phase $+i$ of the G3 coefficients. All quoted coefficients are the machine-checked ones.
\end{remark}

\begin{remark}[Placement sensitivity of the G2 slot map]\label{rem:ss:ab:slotmap}
The decomposition of the G2 word $\tfrac13\beta_r(4p-q)^{\dot\alpha}\partial_{\dot\alpha}c$ is placement-sensitive: on the graph-word placement $(\partial^{\dot\alpha}\beta_r)\partial_{\dot\alpha}c$ the pair coefficient is $-\tfrac13$, while on the canonical pairing $\langle\beta_r,\partial c\rangle=-(\partial^{\dot\alpha}\beta_r)\partial_{\dot\alpha}c$ of~\eqref{eq:ss:ab:paireq} it is $+\tfrac13$, and the consistent md-level total-derivative quotient of the displayed word on the canonical carrier is $\tfrac53$. The intermediate slot-map displays mix the two placements (the typed slot map of the Omega-routing audit versus the identity $(P^{\dot a}B_1)D_{\dot a}=-\langle B_1,D\rangle$ of the external-divergence audit). The displayed G2 raw pair/EOM entries and the quotient value $+1$ in~\eqref{eq:ss:ab:flowtd} are therefore machine-sealed by the typed replay rather than derived at md level. The renormalized endpoint does not depend on this link: the Ward ray and its unit scale are anchored by the descendant identity~\eqref{eq:ss:triplet}, and the full renormalized vector matches the holomorphic-twist target word by word in the 81/81 machine round-trip.
\end{remark}

\paragraph{8. The final result.} Restoring the letter normalizations~\eqref{eq:ss:letters} --- the insertions contribute $(-i/\sqrt2)(1/\sqrt2)=-i/2$, each dotted output pair contributes $-i\sqrt2$, and the $\gamma$ pairs contribute $1$, so the unit-normalized vector is $-i\sqrt2$ times the bracket and $(-i/2)(-i\sqrt2)=-1/\sqrt2$ --- the renormalized vector~\eqref{eq:ss:ab:flowren} becomes the channel-table row
\begin{equation}\label{eq:ss:ab:result}
\boxed{\;
\nabla_{-}\big(b^{A}\beta_r^{B}\big),\quad
\nabla_{-}\big(\beta_r^{A}b^{B}\big)
=-\frac{\h g^{2}}{16\sqrt2\,\pi^{2}}\,f^{ACD}f^{BCE}\,
\big\{\langle\beta_{r},\partial c\rangle
+\langle\partial c,\beta_{r}\rangle
+\varepsilon_{rst}\langle\gamma^{s},\gamma^{t}\rangle\big\}^{DE} ,\;}
\end{equation}
in exact agreement with~\eqref{eq:ss:ch-bbeta}. The two orderings are equal: the source-order Koszul sign is $(-1)^{|b||\beta|}=(-1)^{0\cdot1}=+1$, the superpotential flavor/color reflection contributes $+1$, the simplex moments are invariant under the reflection relabeling, and the reflected vector maps word by word onto~\eqref{eq:ss:ab:flowren}.

\paragraph{9. Comparison with the holomorphic twist.} Budzik and Kulp print, for this family, \cite[eq.~(3.68)]{BK}
\begin{equation}\label{eq:ss:ab:bk}
Q_1\big((\beta_I)^{A}b^{B}\big)
=\kappa^{2}f_{ACD}f_{BCE}
\Big[\partial_{\dot\alpha}(\beta_I)^{D}\partial^{\dot\alpha}c^{E}
+\partial_{\dot\alpha}c^{D}\partial^{\dot\alpha}(\beta_I)^{E}
+\varepsilon_{IJK}\,\partial_{\dot\alpha}(\gamma^{J})^{D}
\partial^{\dot\alpha}(\gamma^{K})^{E}\Big] .
\end{equation}
Under the letter dictionary of Section~\ref{sec:ht}, the three words of~\eqref{eq:ss:ab:bk} map exactly onto $\langle\beta_r,\partial c\rangle+\langle\partial c,\beta_r\rangle$ and $\varepsilon_{rst}\langle\gamma^{s},\gamma^{t}\rangle$, and BK's prefactor $\kappa^{2}$ maps onto $-\h g^{2}/(16\sqrt2\,\pi^{2})$: the match is word by word, for both orderings. The reversed ordering $Q_1(b^{A}(\beta_I)^{B})$ is not printed as a final display in \cite{BK}; it is fixed there by the descent relations, and the machine round-trip confirms the identical canonical word for both directions. Two source inconsistencies of \cite{BK} are recorded in Section~\ref{sec:ht}: the missing $\Gamma(3)=2$ in their printed triangle kernel, and the $-\tfrac14$ versus $-\kappa^{2}$ clash, adjudicated to the main branch $-\kappa^{2}$.

\begin{remark}[Status of the comparison]\label{rem:ss:ab:htstatus}
The match with~\eqref{eq:ss:ab:bk} is a cross-check against an external target, not a completeness proof of the graph census: the adversarial seventeen-candidate census review remains open, and no holomorphic-twist coefficient enters the derivation of~\eqref{eq:ss:ab:result}. Before the finite shift~\eqref{eq:ss:ab:flowshift}, the raw $G_{3,2},G_{3,3}$ coefficients are exactly twice their twist counterparts in every common basis (since $\varepsilon_{1JK}X_{JK}=X_{23}-X_{32}$); the shift is what removes this factor two.
\end{remark}
